\documentclass[11pt]{article}
\usepackage{siunitx}
\usepackage{cleveref}
\usepackage[margin=1in]{geometry}
\usepackage{graphicx}
\usepackage{booktabs}
\usepackage{float}
\usepackage{caption}
\usepackage{subcaption}

% Change only these lines to point the note at another rendered report.
\input{hundredpatches_50iters_paths.tex}

\title{Rainfall Estimation}
\author{}
\date{}

\begin{document}
\maketitle

This note summarizes results for a collection of 100 patches. All aggregate figures and tables that follow should therefore be read as patch-level summaries over this hundred-patch evaluation set, rather than as statistics from a single patch or from the full archive of available patches.

\section*{Relative Rainy-Pixel Distance Profile}
\begin{figure}[H]
  \centering
  \includegraphics[width=\textwidth]{\distanceboxplot}
  \caption{Rainy-pixel distance-profile box-and-whisker summary for distance to the third closest link. Each patch contributes one median rainy-pixel relative-absolute-error (RAE) value per distance bin, and these patch-level medians are then summarized across patches with a box-and-whisker plot. Here, a rainy pixel is a pixel whose ground-truth rainfall is at least the configured wet threshold (set to 0.6), and relative absolute error for a rainy pixel $p$ is defined as $|GT(p)-PRED(p)| / GT(p)$.}
  \label{fig:RAE}
\end{figure}

This plot summarizes how solver performance changes as pixels get farther from the third-nearest link. Lower values indicate lower rainy-pixel relative absolute error in that distance bin.

\subsection*{Rainy-pixel Count Table for Distance to the 3rd Closest Link}
The table below gives the average number of rainy pixels per patch and the corresponding standard deviation for each distance bin to the 3rd closest link. These are the same patch-level counts summarized in the tick labels of the rainy-pixel distance-profile plots.

\begin{table}[H]
  \centering
    \begin{tabular}{l
                  l
                  S[table-format=5.1]}
    \toprule
    {Distance to third closest link (m)} & {Mean rainy-pixel count per patch (\% of total)} & {SD} \\
    \midrule
    $[0,125]$       & 4267.0 (1.79\%) &   960.4 \\
    $(125,375]$     & 16737.7 (7.01\%) &  3803.2 \\
    $(375,750]$     & 26809.9 (11.22\%) &  6707.0 \\
    $(750,1500]$    & 51988.5 (21.76\%) & 14468.0 \\
    $(1500,3125]$   & 88839.4 (37.18\%) & 27531.1 \\
    $(3125,6000]$   & 45929.4 (19.22\%) & 15679.4 \\
    $(6000,9000]$   & 3835.8 (1.61\%) &  2725.7 \\
    $(9000,\infty)$ & 514.4 (0.22\%) &   966.9 \\
    \midrule
    {Total}         & 238922.0 (100\%) & {} \\
    \bottomrule
  \end{tabular}
  \caption{Distribution of the number of rainy pixels per distance bin. In each patch, rainy pixels are classified into bins according to their distance to the 3rd-closest link. Formally, a pixel $p$ belongs to bin $(d_1,d_2]$ if the disc of radius $d_2$ centered at $p$ intersects at least $3$ links and the disc of radius $d_1$ centered at $p$ intersects at most $2$ links. Let $|\mathrm{bin}_i(P)|$ denote the number of rainy pixels in patch $P$ that belong to $\mathrm{bin}_i$. The table reports, for each bin, the mean and standard deviation of $|\mathrm{bin}_i(P)|$ across patches, with the percentage of the total mean rainy-pixel count shown in parentheses next to the mean. The total is the sum of the per-bin mean counts. The last bin, $(9000,\infty)$ is not shown in \Cref{fig:RAE} because it is only sparsely populated.}
\end{table}

\section*{Confusion Map}
\begin{table}[H]
  \centering
  \begin{minipage}[t]{0.31\textwidth}
    \centering
    \textbf{IDW}

    \vspace{0.35em}
    \scriptsize
    \begin{tabular}{lcc}
      \toprule
       & \multicolumn{2}{c}{Prediction} \\
      GT & Wet & Dry \\
      \midrule
      Wet & $0.893 \pm 0.013$ & $0.002 \pm 0.000$ \\
      Dry & $0.040 \pm 0.004$ & $0.065 \pm 0.010$ \\
      \bottomrule
    \end{tabular}
  \end{minipage}
  \hfill
  \begin{minipage}[t]{0.31\textwidth}
    \centering
    \textbf{ILDW}

    \vspace{0.35em}
    \scriptsize
    \begin{tabular}{lcc}
      \toprule
       & \multicolumn{2}{c}{Prediction} \\
      GT & Wet & Dry \\
      \midrule
      Wet & $0.892 \pm 0.013$ & $0.003 \pm 0.000$ \\
      Dry & $0.037 \pm 0.004$ & $0.068 \pm 0.010$ \\
      \bottomrule
    \end{tabular}
  \end{minipage}
  \hfill
  \begin{minipage}[t]{0.31\textwidth}
    \centering
    \textbf{Nonlinear Optimizer}

    \vspace{0.35em}
    \scriptsize
    \begin{tabular}{lcc}
      \toprule
       & \multicolumn{2}{c}{Prediction} \\
      GT & Wet & Dry \\
      \midrule
      Wet & $0.892 \pm 0.013$ & $0.003 \pm 0.000$ \\
      Dry & $0.036 \pm 0.004$ & $0.069 \pm 0.010$ \\
      \bottomrule
    \end{tabular}
  \end{minipage}
  \caption{Patch-averaged wet/dry confusion matrices for IDW, ILDW, and the nonlinear optimizer at a threshold of $0.6$ mm/h, where a pixel is called wet if its rainfall is at least $0.6$ mm/h and dry otherwise. For each patch, pixels are first assigned to the four confusion categories defined by the ground-truth and predicted wet/dry labels: top-left is TP, top-right is FN, bottom-left is FP, and bottom-right is TN. The count in each category is then divided by the total number of pixels in that patch, producing a patch-level pixel fraction. The table entries report the mean of those patch-level fractions across the 100 evaluated patches, together with the standard error of the mean (SEM). These are therefore fractions of all pixels in a patch, not row-normalized conditional rates.}
\end{table}

\section*{Patch Geometry Overview}
\begin{table}[H]
  \centering
  \begin{tabular}{lrrrrr}
    \toprule
    Approx. patch size (km) & \# links & 1-multiplicity & 2-multiplicity & 3-multiplicity \\
    \midrule
    $50 \times 50$ & 939 & 99 & 840 & 0 \\
    $62 \times 60$ & 1254 & 128 & 1126 & 0 \\
    $72 \times 70$ & 1507 & 157 & 1350 & 0 \\
    \bottomrule
  \end{tabular}
  \caption{Representative link multiplicities for three patches with dimensions close to $50\times 50$, $60\times 60$, and $70\times 70$ km. Links are grouped by endpoint coincidence. A 1-multiplicity link has no coincident partner, a 2-multiplicity link belongs to a duplicated pair, and a 3-multiplicity link belongs to a triplet of coincident links. The table reports numbers of links, not numbers of multiplicity groups.}
\end{table}

\section*{$J_{\mathrm{atten}}$ Ratios}
\begin{table}[H]
  \centering
  \begin{tabular}{lcc}
    \toprule
    Solver & Mean $J_{\mathrm{atten}}$ & SD \\
    \midrule
    IDW & 0.01522 & 0.03808 \\
    ILDW & 0.00289 & 0.00757 \\
    Nonlinear Optimizer & 0.00020 & 0.00079 \\
    \bottomrule
  \end{tabular}
  \caption{Patch-level summary of $J_{\mathrm{atten}}$ across the 100 evaluated patches. For a given patch, $J_{\mathrm{atten}}$ is the attenuation-fit term $J_{\mathrm{atten}} = \frac{1}{N_{\mathrm{valid\ links}}}\sum_{k \in \mathrm{valid\ links}} \frac{\left(\hat{A}_k - A_k^{\mathrm{obs}}\right)^2}{L_k^{\mathrm{km}}}$, where $\hat{A}_k$ is the attenuation implied by the estimated rainfall field on link $k$, $A_k^{\mathrm{obs}}$ is the observed attenuation, and $L_k^{\mathrm{km}}$ is the link length in kilometers. The table reports the mean and standard deviation of this per-patch quantity for each solver. Lower values indicate better agreement with the observed link attenuations.}
\end{table}

\begin{table}[H]
  \centering
  \begin{tabular}{lc}
    \toprule
    Ratio of solver means & Value \\
    \midrule
    Mean($J_{\mathrm{atten}}$(IDW)) / Mean($J_{\mathrm{atten}}$(ILDW)) & 5.28 \\
    Mean($J_{\mathrm{atten}}$(ILDW)) / Mean($J_{\mathrm{atten}}$(Nonlinear Optimizer)) & 14.54 \\
    \bottomrule
  \end{tabular}
  \caption{Ratios formed after first averaging $J_{\mathrm{atten}}$ over patches within each solver. These values therefore compare solver-level means, in contrast to the table below, which averages the per-patch solver-to-solver ratios.}
\end{table}

\begin{table}[H]
  \centering
  \begin{tabular}{lc}
    \toprule
    Ratio & Mean over patches \\
    \midrule
    IDW / ILDW & 16.07 \\
    ILDW / Nonlinear Optimizer & 75.62 \\
    \bottomrule
  \end{tabular}
  \caption{Patch-level ratios of $J_{\mathrm{atten}}$ across the 100 evaluated patches. For a given patch, $J_{\mathrm{atten}}$ is the attenuation-fit term $J_{\mathrm{atten}} = \frac{1}{N_{\mathrm{valid\ links}}}\sum_{k \in \mathrm{valid\ links}} \frac{\left(\hat{A}_k - A_k^{\mathrm{obs}}\right)^2}{L_k^{\mathrm{km}}}$, where $\hat{A}_k$ is the attenuation implied by the estimated rainfall field on link $k$, $A_k^{\mathrm{obs}}$ is the observed attenuation, and $L_k^{\mathrm{km}}$ is the link length in kilometers. For each patch, we compute the solver-to-solver ratio shown in the left column, and the table reports the average of those per-patch ratios across the full hundred-patch set. Because lower $J_{\mathrm{atten}}$ is better, a large value of, for example, ILDW / Nonlinear Optimizer means that the nonlinear optimizer achieves much smaller attenuation-fit error than ILDW on a typical patch.}
\end{table}

\section*{Absolute Attenuation Error per km}
\begin{table}[H]
  \centering
  \begin{tabular}{lcc}
    \toprule
    Solver & Mean & SD \\
    \midrule
    IDW & 0.05034 & 0.03396 \\
    ILDW & 0.01356 & 0.00943 \\
    Nonlinear Optimizer & 0.00273 & 0.00411 \\
    \bottomrule
  \end{tabular}
  \caption{Patch-level summary of the mean absolute attenuation error per km. For a given patch, we compute $\frac{1}{N_{\mathrm{valid\ links}}}\sum_{k \in \mathrm{valid\ links}} \frac{\left|\hat{A}_k - A_k^{\mathrm{obs}}\right|}{L_k^{\mathrm{km}}}$, where $\hat{A}_k$ is the attenuation implied by the estimated rainfall field on link $k$, $A_k^{\mathrm{obs}}$ is the observed attenuation, and $L_k^{\mathrm{km}}$ is the link length in kilometers. Thus, each link contributes equally after its own absolute attenuation error is normalized by link length. The table reports the mean and standard deviation of this per-patch quantity across the 100 evaluated patches for each solver. Lower values indicate better agreement with the observed link attenuations.}
\end{table}

\begin{table}[H]
  \centering
  \begin{tabular}{lcc}
    \toprule
    Solver & Mean & SD \\
    \midrule
    IDW & 0.05512 & 0.04105 \\
    ILDW & 0.00877 & 0.00696 \\
    Nonlinear Optimizer & 0.00252 & 0.00371 \\
    \bottomrule
  \end{tabular}
  \caption{Patch-level summary of the globally length-normalized absolute attenuation error. For a given patch, we compute $\frac{\sum_{k \in \mathrm{valid\ links}} \left|\hat{A}_k - A_k^{\mathrm{obs}}\right|}{\sum_{k \in \mathrm{valid\ links}} L_k^{\mathrm{km}}}$ using the same notation as above. Unlike the previous table, this metric first sums the absolute attenuation error over all links in the patch and then divides by the total link length, so longer links receive more weight. The table reports the mean and standard deviation of this per-patch quantity across the 100 evaluated patches for each solver. Lower values again indicate better agreement with the observed link attenuations.}
\end{table}

\end{document}
